\section{Preliminaries and Motivation}
\label{sec:motivation}

\subsection{Quantizing a WAM}
\label{sec:prelim}

\paragraph{World action models.}
A WAM maps an observation $\bo$ and an instruction $\ell$ to a predicted future latent $\bz$ and an action chunk $\ba$, and the action is conditioned on $\bz$:
\begin{equation}
  \bz = G_{\bW}(\bo, \ell), \qquad \ba = R_{\bW}(\bo, \ell, \bz),
  \label{eq:wam}
\end{equation}
where $\bW$ collects all weights.
We call $R$ the \emph{reader}: the computation that turns the predicted future into an action.
Whether $\bz$ and $\ba$ are generated one after the other or denoised jointly, $R$ is well defined (\cref{sec:exp:setup}).
We write $\bz^\star, \ba^\star$ for the outputs of the full-precision model and $\hat G, \hat R, \hat\bz, \hat\ba$ for those of a quantized one.

\paragraph{Weight-only vector quantization.}
Post-training quantizers round each linear layer $\bW$ to $\hat\bW$ by minimizing a reconstruction error weighted by input statistics,
\begin{equation}
  \min_{\hat\bW}\; \tr\!\big((\hat\bW-\bW)\,\bH\,(\hat\bW-\bW)^\top\big), \qquad \bH = \mathbb{E}\big[\bx\bx^\top\big],
  \label{eq:layer}
\end{equation}
where $\bx$ is the layer input on calibration data~\cite{frantar2023gptq,tseng2024quipsharp,tseng2024qtip}.
Trellis-coded quantization~\cite{tseng2024qtip} stores each layer as
\begin{equation}
  \hat\bW = \operatorname{diag}(\mathbf{s}_v)\,\Phi_m^\top\,\big(\gamma\, C[\mathbf{c}]\big)\,\Phi_n\,\operatorname{diag}(\mathbf{s}_u),
  \label{eq:qtip}
\end{equation}
where $\Phi_m, \Phi_n$ are random Hadamard transforms, $\mathbf{s}_u, \mathbf{s}_v$ are column and row scales, $\gamma$ is a layer scale, and $C[\mathbf{c}]$ decodes the weights from trellis codes $\mathbf{c}$ through a small lookup table $C$ that is shared across layers.
Each code spends $b$ bits per weight.
We report bits per weight (bpw) as the total storage of codes, scales, tables and any full-precision layers divided by the number of weights.

\paragraph{Errors.}
We measure the action error $A = \|\hat\ba - \ba^\star\|^2$ on the executed part of the chunk, the future error $Z = \|\hat\bz - \bz^\star\|^2 / \|\bz^\star\|^2$ on the predicted frames, and the reader error
\begin{equation}
  K = \mathrm{KL}\big(\pi(R(\bz^\star)) \,\big\|\, \pi(R(\hat\bz))\big),
  \label{eq:reader-kl}
\end{equation}
where $R$ is the full-precision reader and $\pi$ maps each standardized action coordinate to a soft histogram.
$K$ measures how much the future error alone changes the action.

\paragraph{Two error paths.}
A weight error can change the action directly, or by degrading the predicted future on which the action is conditioned.
We measure the second path with the \emph{future share}
\begin{equation}
  \rho = 1 - \frac{\mathbb{E}\,\|\hat R(\bz^\star) - \ba^\star\|^2}{\mathbb{E}\,\|\hat R(\hat\bz) - \ba^\star\|^2},
  \label{eq:share}
\end{equation}
the fraction of the action error that disappears when the quantized model reads the full-precision future instead of its own.
It only requires the reader, so it can be computed for any WAM.
Existing quantization objectives measure layer outputs, as in \cref{eq:layer}, or the final action, but never the predicted future on which the action depends.
We next ask what this omission costs.

\subsection{Observations}
\label{sec:observations}

We examine how ultra-low-bit weight quantization affects the predicted future and, through it, the action.
We quantize Fast-WAM~\cite{yuan2026fastwam} with three 2-bit VQ methods~\cite{vanbaalen2024gptvq,tseng2024quipsharp,tseng2024qtip} and measure $A$, $Z$ and $\rho$ on held-out LIBERO contexts.
To locate where future error originates, we additionally raise one module group at a time from 2 to 3 bits.
\cref{sec:supp:prelim} provides the full protocol and results.

\begin{figure*}[t]
  \centering
  \begin{subfigure}[t]{0.32\linewidth}
    \figbox{Placeholder: future share $\rho$ for GPTVQ, QuIP\#, QTIP at 2 and 3 bits (bar chart).}{3.6cm}
    \caption{Restoring the future removes most of the action error.}
    \label{fig:obs:share}
  \end{subfigure}\hfill
  \begin{subfigure}[t]{0.32\linewidth}
    \figbox{Placeholder: gain of raising each module group from 2 to 3 bits, stacked by $A$, $Z$, $K$; video expert $|$ action expert $|$ interface.}{3.6cm}
    \caption{Future error comes from the modules that generate it.}
    \label{fig:obs:modules}
  \end{subfigure}\hfill
  \begin{subfigure}[t]{0.32\linewidth}
    \figbox{Placeholder: diagonal of the reader metric $\bM$ on the first predicted frame, overlaid on both camera views.}{3.6cm}
    \caption{The reader is sensitive to a narrow part of the future.}
    \label{fig:obs:gaze}
  \end{subfigure}
  \caption{What an ultra-low-bit WAM loses (Fast-WAM).
  (a)~Replacing the predicted future of a quantized model with the full-precision one removes most of its action error.
  (b)~Raising the modules that generate the future recovers most of the gain of an extra bit; the modules that only read it leave the future unchanged.
  (c)~The reader is sensitive to a small region of the predicted future.}
  \label{fig:observations}
\end{figure*}

\paragraph{(a) Most action error comes from a degraded future.}
Restoring the predicted future removes 73\% of the action error of 2-bit QTIP and 85\% of that of 3-bit QTIP (\cref{fig:obs:share}).
Across the three VQ methods, the future share at 2 bits ranges from \XX\% to \XX\%.
The reader's own quantization error accounts for the remainder.
This implies that the quantization objective must include the predicted future itself, not only the action.

\paragraph{(b) Future error originates in the modules that generate it.}
Raising the video expert, which generates the future, from 2 to 3 bits recovers 71\% of the reduction in the combined error $A + \lambda Z + \beta K$ (\cref{eq:objective}) that an extra bit buys for the whole model, mostly through $Z$ and $K$ (\cref{fig:obs:modules}).
Raising all 30 blocks of the action expert, which only reads the future, recovers almost nothing and leaves $Z$ and $K$ exactly unchanged.
Conversely, lowering the action expert by one bit leaves the predicted future unchanged.
This implies that the bit budget should follow each module's impact on future fidelity.

\paragraph{(c) The reader is sensitive to a narrow part of the future.}
We linearize the reader at the full-precision future and compute its Gauss--Newton metric $\bM$ (\cref{sec:method:metric}).
In a 37{,}632-dimensional latent, $\bM$ has an effective rank of 41.5; the first predicted frame carries 98\% of its trace, and the largest 10\% of its diagonal entries carry 62\%.
Its diagonal is stable across disjoint halves of the calibration contexts (correlation 0.86), and it concentrates on a small region of the scene (\cref{fig:obs:gaze}).
This implies that future error should be measured in the metric through which the reader sees it, rather than by a uniform error over all latent dimensions.

Together, these observations call for a quantizer that treats the predicted future, as the reader sees it, as a quantization target in its own right.
